\documentclass[aspectratio=169, 10pt]{beamer}

% ---------- Tema y colores ----------
\usetheme{default}
\usecolortheme{default}
\setbeamertemplate{navigation symbols}{}
\setbeamertemplate{footline}[frame number]
\setbeamertemplate{frametitle continuation}{}

\definecolor{StanfordRed}{RGB}{140,21,21}
\definecolor{StanfordGray}{RGB}{77,77,77}
\definecolor{LightGray}{RGB}{240,240,240}
\definecolor{AccentBlue}{RGB}{0,84,166}
\definecolor{AccentGreen}{RGB}{0,130,100}
\definecolor{UCMBlue}{RGB}{4, 171, 226}
\definecolor{UCMNavy}{RGB}{20, 65, 134}

\setbeamercolor{frametitle}{fg=white, bg=UCMBlue}
\setbeamercolor{title}{fg=white}
\setbeamercolor{author}{fg=white}
\setbeamercolor{institute}{fg=white!80}
\setbeamercolor{structure}{fg=UCMBlue}
\setbeamercolor{block title}{bg=UCMBlue!15!white, fg=UCMBlue}
\setbeamercolor{block body}{bg=LightGray}
\setbeamercolor{alerted text}{fg=UCMBlue}
\setbeamercolor{example text}{fg=UCMBlue}

\setbeamerfont{frametitle}{series=\bfseries, size=\large}
\setbeamerfont{title}{series=\bfseries, size=\LARGE}
\setbeamerfont{block title}{series=\bfseries}

% ---------- Paquetes ----------
\usepackage[utf8]{inputenc}
\usepackage[T1]{fontenc}
\usepackage[provide=*,spanish]{babel}

\usepackage{amsmath, amssymb, bm}
\usepackage{tikz}
\usepackage{pgfplots}
\pgfplotsset{compat=1.18}
\usetikzlibrary{arrows.meta, shapes.geometric, positioning, fit, calc,
                decorations.pathreplacing, backgrounds, shadows, patterns}
\usepackage{booktabs}
\usepackage{xcolor}
\usepackage{multicol}
\usepackage{tcolorbox}
\tcbuselibrary{skins, breakable}
% ---------- Comandos útiles ----------
\newcommand{\R}{\mathbb{R}}
\newcommand{\E}{\mathbb{E}}
\newcommand{\Var}{\operatorname{Var}}
\newcommand{\MSE}{\operatorname{MSE}}
\newcommand{\Bias}{\operatorname{Sesgo}}
\newcommand{\y}{\bm{y}}
\newcommand{\X}{\bm{X}}
\newcommand{\bbeta}{\bm{\beta}}

\newcommand{\formula}[1]{%
  \begin{tcolorbox}[colback=UCMBlue!8, colframe=UCMBlue, arc=4pt,
                    boxsep=2pt, left=6pt, right=6pt, top=3pt, bottom=3pt]
  #1
  \end{tcolorbox}
}

%\logo{\includegraphics{figures/logotipo_jpg.jpg}}
% ---------- Portada ----------
\title{\textbf{Machine Learning}}
\subtitle{MDS-211 Magister en Data Science}
\author{Sergio Hernández\\ shernandez@ucm.cl}
\institute{ \begin{tikzpicture}
    \clip [rounded corners=1cm] (0,0) rectangle (4,2);
    \node[anchor=south west, inner sep=0pt] at (0,0) 
        {\includegraphics[width=4cm, height=2cm]{figures/mindslab.png}};
\end{tikzpicture}
\\ \textbf{Universidad Católica del Maule}}
\date{}


\begin{document}
% --- Título ---
{
	\setbeamertemplate{background}{
		\begin{tikzpicture}[remember picture, overlay]
			\fill[UCMBlue] (current page.south west) rectangle (current page.north east);
			\fill[UCMBlue!70!black, opacity=0.5]
			(current page.south west) -- ++(0,3) -- ++(18,0) -- cycle;
			\foreach \i in {1,...,40}{
				\fill[white, opacity=0.04]
				({random()*17}, {random()*10}) circle ({random()*0.3});
			}
		\end{tikzpicture}
	}
	\begin{frame}[plain]
		\vspace{1.2cm}
		\titlepage
		\begin{tikzpicture}[remember picture, overlay]
			\node[white, opacity=0.15, font=\fontsize{80}{80}\selectfont] 
			at (current page.east) [xshift=-2cm] {$\Sigma$};
		\end{tikzpicture}
	\end{frame}
}

\section{Feature Engineering}

\begin{frame}
\frametitle{Qu\'e es Feature Engineering?}
\begin{block}{Definici\'on}
Es el proceso de crear nuevos descriptores a partir de los datos en bruto, de manera de aumentar el poder predictivo de los algoritmos de miner\'ia de datos.
\end{block}

\pause
\begin{itemize}
\item Datos num\'ericos : representan mediciones cuantitativas, por ejemplo cantidad de casos (valores discretos) o registro de temperaturas (valores continuos). 
\pause
\item Datos categ\'oricos : representan valores cualitativos, por ejemplo sexo, estado civil, etc. 
\pause
\item Datos ordinales : contiene datos categ\'oricos y num\'ericos, por ejemplo valores de cuestionarios (respuesta 1,2,3,..).

\end{itemize}
\end{frame}

\section{Datos num\'ericos}
\begin{frame}
\frametitle{Transformaci\'on Logaritmica}
\begin{itemize}
	\item Es una transformaci\'on mon\'otona ya que cambia la escala pero no la posici\'on relativa del dato. 
	\pause
	\item La funci\'on logaritmo comprime el rango de valores altos y expande el rango de valores bajos, por lo tanto permite reducir el sesgo de los datos en bruto. 
	\begin{align}
	x= \operatorname{log}_{10}(x)
	\end{align}
\end{itemize}

\begin{figure}
	\includegraphics[width=0.475\linewidth]{figures/skewed_data}
	\hfill
	\includegraphics[width=0.475\textwidth]{figures/skewed_data_transform}
\end{figure}
\end{frame}

\begin{frame}
\frametitle{Normalizaci\'on}
\begin{itemize}
	\item Los modelos basados en expansiones lineales de las entradas son altamente susceptibles al rango de las variables. Algunas transformaciones comunmente usadas son:
	\begin{itemize}
		\item \textbf{Escalamiento Min-Max} : Comprime los valores originales al intervalo $[0,1]$.
		
		\begin{align}
		x=\frac{x - \operatorname{min}(x)}{\operatorname{max}(x) - \operatorname{min}(x)}
		\end{align}
		\item \textbf{Estandarizaci\'on} : Re-escala los valores originales de modo que tengan media $0$ y  varianza $1$. 
		
		\begin{align}
		x=\frac{x - \operatorname{mean}(x)}{ \sqrt{\operatorname{var}(x)}}
		\end{align}
	\end{itemize}
	 
\end{itemize}

\end{frame}

\begin{frame}
\frametitle{Interacciones}
\begin{itemize}
	\item Los descriptores de interacci\'on calculan el producto de dos descriptores.
	\pause
	\item Por ejemplo si consideramos el siguiente modelo lineal:
	\begin{align}
	y=w_0+w_1 x_1 + w_2 x_2 + w_3 x_3 
	\end{align}
	\pause
	\item Es posible crear nuevos descriptores de interacci\'on, combinando las variables de entrada:
	\begin{align}
	y=w_0+w_1 x_1 + w_2 x_2 + w_3 x_3 + w_{1,2} x_1 x_2 + w_{1,3} x_1 x_3 + w_{2,3} x_2 x_3
	\end{align}	
	
\end{itemize}

\end{frame}

\begin{frame}
\frametitle{Polinomios}
\begin{itemize}
	\item Los descriptores polinomiales calculan polinomios de orden $p$.
	\pause
	\item Por ejemplo si consideramos el siguiente modelo lineal:
	\begin{align}
	y=w_0+w_1 x_1 + w_2 x_2 + w_3 x_3 
	\end{align}
	\pause
	\item Es posible crear nuevos descriptores de interacci\'on, combinando las variables de entrada:
	\begin{align}
	y=w_0+w_1 x_1 + w_2 x_2 + w_3 x_3 + w_{1,1} x_1^2 + w_{2,2} x_2^2 + w_{3,3} x_3^2
	\end{align}	
	
\end{itemize}

\end{frame}

\section{Datos categ\'oricos}
\begin{frame}
\frametitle{One Hot Encoding}
\begin{itemize}
	\item One Hot Encoding es una transformaci\'on para variables categoricas que utiliza $1$ bit para cada categor\'ia. 
	\pause
	\item Si la variable original contiene $k$ categro\'ias, entonces OHE crea $k$ columnas.
	\pause
	\item Las $k$ columnas son dependientes, por lo tanto generalmente se usan las $k-1$ primeras columnas para crear modelos lineales.  
\end{itemize}

\begin{columns}[t]
	\begin{column}{0.3\textwidth}
		\begin{table}[]
			\begin{tabular}{l|l|l}
				Id & Ciudad & Renta \\
				\hline
				1  & Santiago & 2.5\\
				2  & Rancagua & 1.3\\
				3  & Talca   & 1.2 \\
				4  & Santiago & 1.9
			\end{tabular}
		\end{table}
	\end{column}
	\begin{column}{0.7\textwidth}  
	\begin{table}[]
	\begin{tabular}{l|l|l|l|l}
		Id & Santiago & Rancagua & Talca & Renta \\
		\hline
		1  & 1 & 0 & 0 & 2.5\\
		2  & 0 & 1 & 0 & 1.3\\
		3  & 0 & 0 & 1 & 1.2\\
		4  & 1 & 0 & 0 & 1.9
	\end{tabular}
\end{table}
	\end{column}
\end{columns}
\end{frame}

\begin{frame}
\frametitle{Automated Feature Engineering}
\begin{figure}
	\centering
	\includegraphics[width=0.5\linewidth]{figures/features_results_comparison}
	\caption{FeatureTools \footnote{\url{https://github.com/Featuretools/Automated-Manual-Comparison}}}
	\label{fig:featuresresultscomparison}
\end{figure}

\end{frame}
\end{document}